\documentclass[11pt]{article}
\PassOptionsToPackage{obeyspaces,spaces,hyphens}{url}
\usepackage[margin=1in]{geometry}
\usepackage[T1]{fontenc}
\usepackage{tgpagella}
\usepackage{parskip}
\usepackage{amsmath,amssymb}
\usepackage{enumitem}
\usepackage{booktabs}
\usepackage{array}
\usepackage{url}
\usepackage[hidelinks]{hyperref}

\setlist[enumerate]{itemsep=5pt,topsep=4pt}
\setlist[itemize]{itemsep=3pt,topsep=4pt}

\newcommand{\Ro}{\mathrm{Ro}}
\newcommand{\pd}[2]{\frac{\partial #1}{\partial #2}}

\begin{document}

\begin{center}
  {\Large\bfseries A spherical version of the axisymmetric shallow-water model}\\[0.4em]
  {\large Plan for \path|spencerahill/axisym-hc-swm|}\\[0.4em]
  Prepared 2026-11-07, and revised the same day to Spencer's answers on the
  setup and his choices of the two open terms. Re-derived on 2026-11-08
  against the model's \path|main|, after the merge of moist V2.
  This copy covers the model work, phases 1 to 4; the science runs that follow
  are planned in the private repository \path|spencerahill/fixed-ro-eq-area|.
  Symbolic and discrete checks, and quoted numbers, from
  \path|scripts/spherical_swm_plan_checks.py| in that repository.
\end{center}

\section{Executive Summary}

The axisymmetric shallow-water model (SWM) of the Hadley cell's upper branch in
\path|spencerahill/axisym-hc-swm|, after Sobel and Schneider (2009, hereafter
SS09), runs on an equatorial beta plane. The SWM's Coriolis parameter is
$\beta y$, and its equations contain no metric terms and no planetary radius, so
they are exactly the small-angle limit of the spherical equations. The SWM
therefore cannot test a theory of the Hadley cell beyond the small-angle
limit. This plan adds spherical geometry to the SWM
as an option, with every beta-plane run left bit-for-bit unchanged.

\paragraph{Where the spherical equations come from.} SS09 adapted their model
from Xian and Miller (2008), who wrote the same two-layer axisymmetric equations
on the sphere, pole to pole, in the coordinate $\sin\varphi$ for latitude
$\varphi$; Adam and Paldor (2009) used the same coordinate. The spherical SWM
takes its dynamics from Xian and Miller and writes SS09's additional terms on
the sphere. The spherical SWM's coordinate is $y = a\sin\varphi$, for planetary
radius $a$, with
$c = \cos\varphi = (1 - y^2/a^2)^{1/2}$. This $y$ is a coordinate, not
meridional distance, which is $a\,d\varphi = dy/c$. In it the Coriolis parameter
is $f = 2\Omega\sin\varphi = (2\Omega/a)\,y$ exactly, for rotation rate $\Omega$.
The constant $2\Omega/a$ is the equatorial value of the physical
$\beta = 2\Omega\cos\varphi/a$, which varies with latitude. The upper-layer
equations are
\begin{align}
  \pd{u}{t} - v\left(f - \pd{(cu)}{y}\right)
    &= -\mathcal{G}\,u\,\pd{(cv)}{y} - \epsilon_u u - \mathcal{S},
    \label{eq:u}\\
  2\pd{v}{t} + \left(f + \frac{u\tan\varphi}{a}\right) u
    &= -\frac{gH}{T_0}\, c\,\pd{T}{y} + \mathcal{K}_v,
    \label{eq:v}\\
  \pd{\theta}{t} + \frac{\delta\Delta_z}{H}\,\pd{(cv)}{y}
    &= \frac{\theta_E - \theta}{\tau},
    \label{eq:theta}
\end{align}
with $\tan\varphi/a = y/(a^2c)$. Here $\mathcal{G}$ is the existing gate on the
vertical advection of zonal momentum, $\epsilon_u$ the Rayleigh drag
coefficient, $\mathcal{S}$ the parameterized eddy momentum flux divergence, and
$\mathcal{K}_v$ the smoothing of $v$. Setting $c = 1$, dropping the metric term
and writing $f = \beta y$ recovers the current equations symbol for symbol. So
the sphere differs from the beta plane through the factors of $c$ in three
derivatives, the metric term, and the domain, which ends at the poles,
$y = \pm a$. Section~\ref{sec:terms} goes through every term.

Two properties make the spherical SWM a test of full-angle theory. The
angular-momentum-conserving (AMC) wind $\Omega a\sin^2\varphi/\cos\varphi$ is an
exact steady state of the inviscid form of Eq.~\eqref{eq:u}. And the steady
form of Eq.~\eqref{eq:v} is the full-angle gradient wind balance. These
properties, the coordinate identities, the
beta-plane limits and the properties claimed in Section~\ref{sec:terms} were
checked symbolically, with 20 of 20 checks passing. The discrete properties
claimed in Section~\ref{sec:disc} were checked numerically against prototype
operators built from the model's own functions, with 9 of 9 checks passing,
both at the commit the plan was first written against and at the merged
\path|main|, commit \path|0dc0627|.

\paragraph{Settled, from Spencer's answers of 2026-11-07.}
\begin{enumerate}
\item \textbf{Base branch:} the model's \path|main|, into which Spencer merged
  moist V2, the latent-heating version of the moist model, on 2026-11-07, as
  commit \path|0dc0627|.
\item \textbf{Formulation:} every term takes its physically correct spherical
  form. Where physics does not fix the form, Section~\ref{sec:terms} says so
  and gives the options.
\item \textbf{Grid:} uniform in $\sin\varphi$, pole to pole. At the production
  grid spacing of 39.4 km this is 324 intervals, 0.40 of the production grid's
  800, with a latitude spacing of 0.35$^\circ$ at the equator widening to
  6.4$^\circ$ at the poles.
\item \textbf{Planet:} Earth's mean radius, $a = 6371$ km, and Earth's
  rotation rate, $\Omega = 7.292115\times10^{-5}$ s$^{-1}$, the value of the
  WGS84 geodetic system. At that rate one rotation relative to the fixed stars
  takes 86,164.1 s; the stars define an inertial frame, as the Coriolis force
  requires.
  The 86,400 s solar day is 236 s longer, because Earth also advances
  0.99$^\circ$ a day along its orbit, in the same sense as it spins, and must
  turn that much further to face the Sun again. puffins'
  \path|ROT_RATE_EARTH| is $2\pi/86{,}400$ s, the solar-day value, 0.27\%
  smaller. Then $2\Omega/a = 2.29\times10^{-11}$ m$^{-1}$ s$^{-1}$, and the
  SWM's default $\beta$ of $2\times10^{-11}$ is 13\% below it, so every
  beta-plane run used for comparison is a new run at $\beta = 2\Omega/a$.
\item \textbf{The eddy momentum flux divergence $\mathcal{S}$}
  (Section~\ref{sec:emfd}): the divergence of an eddy angular-momentum flux,
  the only one of three spherical forms of SS09's empirical closure that
  conserves angular momentum. For a uniform-$\Ro$ wind the three differ by up
  to 17\% at 30$^\circ$.
\item \textbf{The smoothing of $v$} (Section~\ref{sec:v}): the vector
  Laplacian. SS09 added the smoothing ``to reduce numerical noise'', so no
  physical process fixes its form.
\end{enumerate}

\paragraph{Open, for Spencer.} One choice I would make unless Spencer objects: the forcing on the sphere is
the spherical profile that all three of the SWM's forcing profiles were adapted
from (Section~\ref{sec:thermo}), so it needs no decision.

\paragraph{The approach.} The spherical model re-discretizes the current scheme
and keeps every numerical choice in it: leapfrog stepping with the Asselin
filter, the staggered grid with $v$ on cell faces, the monotonized-central (MC)
stencil of the eddy term, first-order upwind meridional advection, and both
Heaviside gates. Each operator is written so that, with $c = 1$ and the metric
term off, it performs the same floating-point operations in the same order as
the current one, apart from multiplications by exactly 1.0, which return their
operand unchanged. A run in that flat configuration must then reproduce, bit
for bit on the same machine, the output of the unmodified model in the four
regression configurations, which checks the restructured code before any
curvature enters. (The stored regression baselines themselves are reproduced
only to cross-machine roundoff, a relative $10^{-11}$, so the bit-for-bit
reference is the unmodified model run beside the new one.) Any departure from the current numerics is
tested separately, one at a time: in the June 2026 rewrite of the numerics,
three untested changes bundled into one plan all had to be reverted.

\paragraph{Phases.}
\begin{enumerate}
\item \textbf{Specification.} The operators term by term (Section~\ref{sec:terms}),
  their discretization including the polar cells (Section~\ref{sec:disc}), and
  a file-by-file change list (Section~\ref{sec:files}). \emph{Verified by} a
  symbolic check of every continuous claim and a numerical check of every
  discrete one. Complete with this revision.
\item \textbf{numpy reference.} Geometry configuration, metric arrays,
  operators, output, restart files and the command line. \emph{Verified by}
  unit tests per operator, among them exact discrete conservation of water and
  exact mirror parity, and by the flat configuration reproducing the
  unmodified model bit for bit on the same machine.
\item \textbf{Validation runs on berimbau}, each with pass criteria set before
  it runs (Section~\ref{sec:valid}): a drag-limited angular-momentum identity
  with the eddy term and vertical advection off; steady budgets; the timestep
  ceiling and a resolution twin; and the scaling of the difference between the
  sphere and a beta-plane twin with the width of the cell. Every run uses
  Earth's $a$ and $\Omega$.
\item \textbf{numba kernel.} The fused day kernel with the spherical operators.
  \emph{Verified by} bitwise parity with the numpy reference, through the
  existing parity suite.
\item \textbf{Science runs,} in a separate research project, which takes over
  once phases 2 to 4 are done (Section~\ref{sec:handoff}).
\item \textbf{Later, optional.} Moist V1 and V2 and the seasonal forcing on the
  sphere, verified as in phases 2 to 4.
\end{enumerate}

\paragraph{Cost.} The spherical grid has 0.40 of the production grid's points,
so a run should cost no more per model day than the current 0.086 s on
berimbau. That cost is an estimate pending phase 3's smoke timing: the metric
arithmetic adds work per point, and the timestep ceiling near the poles is
unmeasured. For development, the closest precedents in this model are the
staggered grid and the numba backend, which took about two days of sessions
each, 2026-07-11 to 07-12 and 2026-07-13 to 07-14. Phases 2 to 4 together are
comparable to both combined, plus the polar cells, which have no precedent
here. That comparison is a precedent only; no part of this work has been
timed.

\section{The spherical equations, term by term}
\label{sec:terms}

For each term, this section gives the current beta-plane form, as written in
SS09 with the correction of Sobel and Schneider (2013, hereafter SS13) and as
implemented in the model; the spherical form; its source; and whether physics
fixes it. Every spherical form follows from four identities in
$y = a\sin\varphi$, each checked symbolically:
\begin{equation}
  \zeta = -\pd{(cu)}{y},\qquad
  \nabla\cdot v = \pd{(cv)}{y},\qquad
  \frac{1}{a}\pd{}{\varphi} = c\,\pd{}{y},\qquad
  \frac{u\tan\varphi}{a} = \frac{y\,u}{a^2 c},
  \label{eq:identities}
\end{equation}
for relative vorticity $\zeta = -(a\cos\varphi)^{-1}\partial_\varphi(u\cos\varphi)$
and divergence $\nabla\cdot v = (a\cos\varphi)^{-1}\partial_\varphi(v\cos\varphi)$.
Grid cells of equal width in $y$ have equal area, so the area-weighted budgets
of mass, heat and water take the same discrete form as on the plane.

\subsection{Zonal momentum}

\paragraph{Coriolis force and meridional advection.} Beta plane:
$-v(\beta y - \partial_y u)$. Sphere: $-v(f + \zeta) = -v[f - \partial_y(cu)]$,
the meridional flux of absolute vorticity. Source: Xian and Miller's
Eqs.~(3) and (4), the advection of angular momentum $M = a\cos\varphi\,(\Omega a\cos\varphi + u)$
by $v$, which is this term written for $M$. \emph{Fixed by physics.} The model
keeps its split into a Coriolis term and first-order upwind advection, the
latter now of $cu$.

\paragraph{Vertical advection of zonal momentum.} Beta plane:
$-\mathcal{G}\,u\,\partial_y v$, where mass entering the upper layer from below
carries zero zonal velocity. Sphere: $-\mathcal{G}\,u\,\partial_y(cv)$. Source:
Xian and Miller's Eq.~(4), in which rising air carries the lower layer's angular
momentum $M_B = \Omega a^2\cos^2\varphi$, that is, zero relative angular
momentum; since $M - M_B = a\cos\varphi\,u$, the term for $u$ is the one above.
\emph{Fixed by physics.} Both gates apply unchanged: Xian and Miller's gate is
SS09's kinematic one, on where the divergence is positive.

\paragraph{Rayleigh drag.} $-\epsilon_u u$ on both, a local term. \emph{Fixed.}

\subsection{The eddy momentum flux divergence}
\label{sec:emfd}

Beta plane: $\mathcal{S} = v_d\,\mathcal{H}(u)\,\mathrm{sgn}(y)\,\partial_y u$
(SS09 Eq.~2.5), which SS09 describe as ``empirical, motivated by the eddy
momentum flux divergence in the idealized GCM integrations of Schneider and
Bordoni (2008)''. \emph{Not fixed by physics.} On the plane the term has two
readings, and on the sphere each reading gives a different form, beside a
third, literal one:
\begin{enumerate}[label=(\alph*)]
\item \textbf{A closure on the eddy vorticity flux.} By the Taylor identity, the
  eddy momentum flux divergence is minus the eddy flux of relative vorticity
  for nondivergent eddies, on the plane and the sphere alike. SS09's term sets that flux to
  $v_d\,\mathcal{H}(u)\,\mathrm{sgn}(y)$ times the mean relative vorticity, and
  the same statement on the sphere gives
  \[
    \mathcal{S} = -v_d\,\mathcal{H}(u)\,\mathrm{sgn}(y)\,\zeta
      = v_d\,\mathcal{H}(u)\,\mathrm{sgn}(y)\,\pd{(cu)}{y}.
  \]
  Form (a) keeps the drag-free steady closure $\Ro = v/(v + v_d)$ exact.
\item \textbf{The literal transcription,} with $\partial_y$ read as a derivative
  along meridional distance:
  \[
    \mathcal{S} = v_d\,\mathcal{H}(u)\,\mathrm{sgn}(y)\,c\,\pd{u}{y}.
  \]
  That closure becomes $\Ro = [v - v_d u/(2\Omega a c)]/(v + v_d)$.
\item \textbf{The divergence of an eddy angular-momentum flux.} SS09's term is,
  away from the equator, the beta-plane divergence of the eddy momentum flux
  $\overline{u'v'} = v_d\,\mathcal{H}(u)\,\mathrm{sgn}(y)\,u$. The spherical
  divergence of that flux is
  \[
    \mathcal{S} = \frac{1}{a\cos^2\varphi}\pd{}{\varphi}
      \left(\cos^2\varphi\;\overline{u'v'}\right)
      = \frac{v_d\,\mathcal{H}(u)\,\mathrm{sgn}(y)}{c}\,\pd{(c^2 u)}{y}.
  \]
  Because $c\,\mathcal{S}$ is then an exact derivative in $y$, the global torque
  it exerts reduces to boundary values, which vanish where $u = 0$ bounds the
  westerlies, just as on the plane. The closure becomes
  $\Ro = [v + v_d u/(2\Omega a c)]/(v + v_d)$.
\end{enumerate}
For a uniform-$\Ro$ wind, forms (b) and (c) equal form (a) times
$1 \pm \tan^2\varphi/2$: a difference of 2\% at 10$^\circ$, 7\% at 20$^\circ$
and 17\% at 30$^\circ$. \emph{Chosen by Spencer, 2026-11-07:} (c). Eddies
redistribute angular momentum and cannot create it, and (c) is the only form
that respects that; the exact closure of (a) is a convenience for theory. The
reading of SS09's term
as a vorticity closure in (a), and as a flux divergence in (c), are mine; SS09
give neither.

\subsection{Meridional momentum}
\label{sec:v}

\paragraph{Tendency, Coriolis, metric and pressure-gradient terms.} Beta plane:
$2\partial_t v + \beta y u = -(gH/T_0)\,\partial_y T$. Sphere:
$2\partial_t v + (f + u\tan\varphi/a)\,u = -(gH/T_0)\,c\,\partial_y T$. Source:
Xian and Miller's Eq.~(9), from integrating the vertical derivative of the
meridional momentum equation over the troposphere, with the lower layer's
meridional wind equal and opposite to the upper layer's and the surface zonal
wind negligible. SS13 add that the factor of 2 requires the two layers to be
equally thick. The advection $v\,\partial v/\partial y$, which SS13 removed,
cancels exactly between the two layers when the lower layer's meridional wind is
minus the upper layer's, on either geometry; that reason is my own derivation. The metric term
$u^2\tan\varphi/a$ comes from the upper layer alone. \emph{Fixed by physics},
given SS09's two-layer assumptions.

\paragraph{The smoothing of $v$.} Beta plane: $k_v\,\partial_y^2 v$, which SS09
added ``to reduce numerical noise''. \emph{Not fixed by physics}, since the term
represents no physical process. Two candidates: (a)~the meridional component
of the vector Laplacian, which for any axisymmetric field is the meridional
gradient of the divergence, since the rotational part $\mathbf{k}\times\nabla\zeta$
is then purely zonal:
$k_v\,c\,\partial_y[\partial_y(cv)] = k_v[\nabla^2 v - v/(a^2\cos^2\varphi)]$;
(b)~the scalar Laplacian applied to $v$ as though it were a scalar, the literal
transcription of $\partial_y^2 v$, which drops the curvature term.
\emph{Chosen by Spencer, 2026-11-07:} (a). The vector Laplacian is a vector operator, so its
form does not depend on the coordinate, and it acts on the divergence, which is where the grid-scale
noise appears. Whether it damps that noise as well as the current form does is
a hypothesis, to be tested against the model's $v$-smoothness monitor and its
timestep ceiling.

\subsection{Thermodynamics}
\label{sec:thermo}

\paragraph{Adiabatic warming and cooling.} Beta plane:
$(\delta\Delta_z/H)\,\partial_y v$. Sphere: $(\delta\Delta_z/H)\,\partial_y(cv)$,
from the spherical continuity equation, Xian and Miller's Eq.~(6).
\emph{Fixed by physics.} SS09's thermodynamic equation has no horizontal
advection of $\theta$, unlike Xian and Miller's Eq.~(12); that is SS09's
modeling choice, independent of the geometry, and it is kept.

\paragraph{Newtonian cooling.} $(\theta_E - \theta)/\tau$ on both, a local term.
\emph{Fixed.}

\paragraph{The forcing $\theta_E$.} The SWM's three profiles are beta-plane
renderings of one spherical profile, that of Schneider and Bordoni's (2008)
idealized GCM, $\hat{T}_{se} - \Delta_h(\sin^2\varphi - 2\sin\varphi_0
\sin\varphi)$ for a heating maximum at $\varphi_0$. The model's seasonal profile,
and its sin$^2$ profile with no offset, are that profile with $\sin\varphi$
replaced by $\sin(\pi y/2y_1)$, so that $y_1$ stands in for the pole. The
parabolic profile of SS09 is its small-angle form,
$\sin^2\varphi \approx (y/a)^2$, with $y_1$ in the place of $a$. On the sphere the three therefore collapse to
$\theta_E = \theta_{00} - \Delta_y(\sin^2\varphi - 2\sin\varphi_0\sin\varphi)$,
with $\Delta_y$ the equator-to-pole drop for $\varphi_0 = 0$, as SS09 define it.
In the coordinate $y = a\sin\varphi$ that profile is
$\theta_{00} - \Delta_y(y^2 - 2y_0 y)/a^2$, with $y_0 = a\sin\varphi_0$: SS09's
parabolic profile with $y_1 = a$, exact on the sphere, and with no plateau,
since $y$ ends at $\pm a$. The half-width
$y_1$ and the parabolic profile's plateau have no other spherical
counterpart. That spherical forcing is the choice I would make unless Spencer
objects.

\paragraph{Eddy heat diffusion.} Beta plane: $\kappa_\theta\,\partial_y^2\theta$.
Sphere: $\nabla\cdot(\kappa_\theta\nabla\theta) =
\partial_y(\kappa_\theta\,c^2\,\partial_y\theta)$, which follows from the
identities~\eqref{eq:identities}. \emph{Fixed}, as the diffusion of a scalar.

\subsection{Moisture, for moist V1 and V2}

Beta plane: $\partial_t W + \partial_y[-(2a_W - 1)\,vW - D\,\partial_y W] = E_0 - P$,
for column water vapor $W$ with lower-layer fraction $a_W$, eddy diffusivity $D$,
evaporation $E_0$ and precipitation $P$. Sphere:
$\partial_t W + \partial_y\{c[-(2a_W - 1)\,vW - D\,c\,\partial_y W]\} = E_0 - P$,
the spherical divergence of the same flux. Precipitation and the latent heating
of moist V2 are local. \emph{Fixed by physics.}

\subsection{The poles}

Xian and Miller (2008) and Adam and Paldor (2009) both run pole to pole, with
the meridional wind zero at the poles. On the staggered grid, the outermost $u$
and $\theta$ points sit at the poles. There $u = 0$ is the regularity condition,
since the relative vorticity contains $u\tan\varphi/a$ and stays finite only if
$u$ vanishes at the poles, and it coincides with the wall condition the model
already imposes. The
meridional wind stays on interior faces, and since $c = 0$ at the poles, no flux
crosses them.

\subsection{Diagnostics}

The local Rossby number becomes $\Ro = -\zeta/f$ with the spherical $\zeta$.
Jet and cell-edge latitudes are
$\arcsin(y/a)$. The output keeps $y$ in meters as its coordinate, so that
existing analysis code reads the files unchanged, and adds latitude beside it.

\section{Discretization on the grid uniform in sin(latitude)}
\label{sec:disc}

Every operator below is transcribed from the model's code at the tip of the
moist V2 branch, commit \path|6503a62| of \path|feature/moist-v2-sweep|, with
the function it modifies named. The merge into \path|main|, commit
\path|0dc0627|, changed none of these functions: its changes to
\path|ss09/sw_model.py| are confined to the moist V2 relaxation target, in
\path|radiative_target_profile| and the model's constructor.

\subsection{The grid and its geometry arrays}

The spherical model keeps the staggered layout of the current one: $u$,
$\theta$ and the column water vapor $W$ at $n_y$ cell centers, $v$ at the
$n_y - 1$ faces between them, and the outermost centers on the boundaries,
which are now the poles, $y = \pm a$. At the production equatorial spacing
this gives $n_y = 325$ centers and $\Delta y = 2a/324 = 39.33$ km. Since the
coordinate is $a\sin\varphi$, $\Delta y$ is the distance between centers only at
the equator. Elsewhere the distance is $\Delta y/\cos\varphi$: 45.4 km at
30$^\circ$ and 78.7 km at 60$^\circ$. The outermost interior centers sit at
83.63$^\circ$, and the outermost faces at 85.50$^\circ$.

The centers are $y_k = a\,(k/m)$ for the integers $k = -m, \dots, m$, with
$m = 162$. The model's configuration builds its grid with \path|np.linspace|,
which is exactly antisymmetric about the equator for the production beta plane
($n_y = 801$) but not on the sphere: \path|np.linspace(-a, a, 325)| puts 248 of
its 325 points off minus their mirror images, by up to $1.9\times10^{-9}$ m.
The model's exact mirror parity, under which a hemispherically symmetric
forcing keeps $u$ and $\theta$ exactly symmetric and $v$ exactly
antisymmetric, requires an exactly antisymmetric $y$, as the comment on
\path|y_v| in \path|sw_config.py| states. The values $a\,(k/m)$ are
antisymmetric by construction, since negating $k$ negates the quotient
exactly, and the endpoints are exactly $\pm a$.

The configuration computes these geometry arrays once:
\begin{itemize}
\item $c = [(1 - y/a)(1 + y/a)]^{1/2}$ at the centers and the faces, exactly
  zero at the pole centers and exactly symmetric;
\item $c^2$ at the centers;
\item $1/c$ at the interior centers, set to zero at the two pole centers, whose
  $u$ tendency the condition $u = 0$ there discards;
\item $f = (2\Omega/a)\,y$ at the centers and the faces;
\item the metric factor $\tan\varphi/a = y/(a^2c)$ at the faces.
\end{itemize}
On the beta plane the same arrays are $c = c^2 = 1/c = 1$ exactly and
$f = \beta y$, and the metric term is skipped. A multiplication by exactly 1.0
returns its operand unchanged in IEEE arithmetic, so every operator below then
reproduces the current one bit for bit. The checks script tests this on
prototype operators built from the model's own functions. On 20 random states
with exact and signed zeros planted, for both gates on the vertical advection,
and with the eddy heat diffusion off and on, every prototype tendency of $u$,
$v$, $\theta$ and $W$ equals the model's bit for bit. Two plausible slips both
fail that test: the smoothing computed in an algebraically equal but
re-associated form, and the factor of $c$ in the pressure gradient applied
after the division by $T_0$.

\subsection{The operators}

Notation, for a face field $g$ and a center field $q$, with the pole centers at
$k = \pm m$ and the faces at half-integer indices:
\begin{itemize}
\item $\mathrm{div}_k[g] = (g_{k+1/2} - g_{k-1/2})/\Delta y$ at the interior
  centers, and $2g_{-m+1/2}/\Delta y$ and $-2g_{m-1/2}/\Delta y$ at the two
  poles (\path|v_divergence_at_centers|). Each pole value is the finite-volume
  divergence of a half cell with one face, and since cells of equal width in
  $y$ have equal area, $\mathrm{div}$ is the finite-volume divergence on the
  sphere.
\item $\bar{g}_k$ and $\bar{q}_{k+1/2}$: the two-point averages between the
  grids (\path|v_faces_to_centers|, and the face average of $u$ in
  \path|dv_dt|), with $\bar v = 0$ at the poles.
\item $\delta^{\mathrm{up}}_k[q]$: the first-order upwind difference of $q$ on
  the sign of $\bar v_k$ (\path|merid_advec_u|).
\item $M_k[q]$: the MUSCL derivative with monotonized-central limited slopes,
  upwinded from the equatorward side (\path|muscl_mc_du_dy|).
\item $L_{k+1/2}[g]$: the three-point face Laplacian in the mirror-symmetric
  association, with mirror ghosts beyond the outermost faces
  (\path|v_face_laplacian|).
\end{itemize}

\paragraph{Coriolis force and meridional advection of $u$} (\path|du_dt|,
\path|merid_advec_u|). Now $\beta y_k\,\bar v_k - \bar v_k\,\delta^{\mathrm{up}}_k[u]$;
on the sphere $f_k\,\bar v_k - \bar v_k\,\delta^{\mathrm{up}}_k[cu]$.

\paragraph{Vertical advection of $u$} (\path|vert_advec_u|).
$\mathcal{G}_k\,u_k\,\mathrm{div}_k[cv]$, with the kinematic gate reading
$\mathrm{div}_k[cv]$.

\paragraph{The eddy term, form (c)} (\path|edd_mom_flux_div_u|).
$\mathcal{S}_k = v_d\,\mathcal{H}(u_k)\,\mathrm{sgn}(y_k)\,c_k^{-1}M_k[c^2u]$.
The MUSCL derivative is a difference of reconstructed face values: in the
north, $M_k[q] = (q^L_{k+1/2} - q^L_{k-1/2})/\Delta y$, with
$q^L_{k+1/2} = q_k + \tfrac12\Delta y\,\sigma_k$ the reconstruction from the
equatorward cell and $\sigma$ the limited slope, and in the south the mirror
form. So $c_k\mathcal{S}_k\,\Delta y$ is a difference of face values within
each hemisphere's westerly region, and the sum over that region equals $v_d$
times its boundary values, as in the continuous form. The checks script
confirms this for $u = U\sin\varphi\cos\varphi$ to a relative
$5\times10^{-17}$, and the literal form (b) fails the same check. The limiter
acts on $c^2u$, the quantity form (c) differentiates, so the existing function
is reused unchanged. The alternative, limiting $u$ and multiplying its face
values by $c^2$, also reduces to the current operator on the plane, and it is
the first change to test if the spherical eddy term proves noisy.

\paragraph{Meridional momentum} (\path|dv_dt|). At each face,
\[
  2\,\partial_t v = -\Bigl(f + \frac{\tan\varphi}{a}\,\bar u\Bigr)\bar u
    - \frac{gH}{T_0}\,c\,\frac{T_{k+1} - T_k}{\Delta y} + k_v\,c\,L[cv].
\]
The smoothing $c\,L[cv]$ is algebraically
$c\,(\mathrm{div}_{k+1}[cv] - \mathrm{div}_k[cv])/\Delta y$, the face gradient
of the discrete divergence, because the mirror ghost reproduces the half-cell
divergence at the pole. The smoothing is therefore the discrete form of the chosen vector
Laplacian, and it is exactly dissipative: summed over the faces, $v$ times
$c\,L[cv]$ equals minus the sum over centers of $\mathrm{div}[cv]^2$, with half
weights at the two pole half cells. So the smoothing never increases the sum
of $v^2$ over the faces, and lowers it whenever the divergence is nonzero. The
checks script confirms both
identities to relative differences of $3\times10^{-16}$ and $1\times10^{-16}$. With $c = 1$ they hold for the
current smoothing too, which is already the gradient of the discrete
divergence.

\paragraph{Thermodynamics.} The adiabatic term, in
\path|vert_advec_theta|, is $-(\delta\Delta_z/H)\,\mathrm{div}_k[cv]$. The eddy
heat diffusion, in \path|eddy_heat_flux| and off by default, becomes
$\kappa_\theta\,\partial_y(c^2\partial_y\theta)$, with the current centered
\path|np.gradient| applied twice.

\paragraph{Moisture} (\path|moisture_transport_tendency|). The face flux is
$F = c\,[\,c_W W^{\mathrm{MC}} - D\,c\,(W_{k+1} - W_k)/\Delta y\,]$, with
$c_W = -(2a_W - 1)v$ the column transport velocity and $W^{\mathrm{MC}}$ the
current limited upwind face value; the tendency is $-\mathrm{div}_k[F]$. Total
water is conserved exactly, since the flux vanishes at the poles, and the
half-cell weights of \path|cwv_integral| are the equal-area weights. The checks
script confirms this to a relative $3\times10^{-18}$.

\paragraph{Local terms.} Rayleigh drag, Newtonian cooling, precipitation and
the latent heating of moist V2 are unchanged.

\paragraph{Forcing} (\path|theta_e.py|). A new profile,
$\theta_{00} - \Delta_y(s^2 - 2s_0s)$ with $s = y/a$ and $s_0 = \sin\varphi_0$.
Moist V2 relaxes to this profile itself, as it relaxes to the current
profiles since the merge; an explicit radiative target, from
\path|--delta-y-rad|, replaces $\Delta_y$ with $\Delta_y^{\mathrm{rad}}$, as
now.

\subsection{The polar cells}

Each pole center sits at the edge of a half cell covering 0.154\% of the
sphere's area, 1/648. There $u = 0$, as at the current walls. No flux crosses a pole, since
$c = 0$ there, so the pole divergence is that of a half cell with one face.
The eddy term is zero at the poles, through $1/c = 0$; the condition $u = 0$
discards the pole tendency anyway, and a zero keeps $u$'s leapfrog history
there finite.

Near a pole, a regular axisymmetric field is even or odd in the colatitude
$\chi$. Since $y/a = \cos\chi$ and $\chi^2$ is an analytic function of
$1 - y/a$, an even field is a smooth function of $y$ up to the pole, while an
odd field is $(a - y)^{1/2}$ times a smooth function. The winds $u$ and $v$ are
odd, and so is $c^2u$; $cu$, $cv$, $\theta$, $T$ and $W$ are even. The
divergence, the pressure gradient, the advection of $cu$, the smoothing and the
moisture fluxes all act on even fields, so the uniform grid in $y$ resolves
them up to the poles. Three operations act on odd fields and lose accuracy in
the outermost cells: the averages of $u$ and $v$ between centers and faces, and
the eddy term's derivative of $c^2u$. For the regular field
$v = V\sin\varphi\cos\varphi$, the face-to-center average of $v$ is off by
3.5\% at 83.6$^\circ$, 0.85\% at 81.0$^\circ$ and 0.22\% at 77.2$^\circ$,
while within 60$^\circ$ of the equator its error stays below $10^{-4}$ of
$\max|v|$. The unit tests of phase 2 measure each operator's convergence order
on regular test fields over the whole domain, polar cells included.

\subsection{The timestep}

In the coordinate $y$, gravity waves travel at $c$ times their physical speed,
because the pressure gradient and the divergence each carry one factor of $c$.
With the spacing in $y$ uniform, the wave limit on the timestep is then set at
the equator and equals the beta plane's at the same $\Delta y$. The advection
speed in $y$ is $cv$, and the smoothing's rate carries $c^2$, so neither grows
toward the poles. The one term that does grow is the metric frequency
$u\tan\varphi/a$: at the outermost faces, where $\tan\varphi = 12.7$, it is
$2.0\times10^{-5}$ s$^{-1}$ for $u = 10$ m s$^{-1}$ and
$6.0\times10^{-5}$ s$^{-1}$ for 30 m s$^{-1}$, 0.14 and 0.41 of $2\Omega$.
These are expectations from the equations: phase 3 finds the ceiling by
bisection, and no run configuration is chosen from them before then.

\section{Changes, file by file}
\label{sec:files}

First written against commit \path|6503a62|, and re-derived on 2026-11-08
against the merged \path|main|, commit \path|0dc0627|, with no change to the
list. Between the two commits, \path|ss09/sw_config.py| changed two moist
defaults, the lower-layer water fraction $a$ from 0.85 to 0.79 and the
critical column water vapor $W_c$ from 50 to 40 kg m$^{-2}$, and made the
radiative contrast $\Delta_y^{\mathrm{rad}}$ optional where it had defaulted
to 75 K; \path|ss09/sw_model.py| changed the moist V2 relaxation target;
\path|ss09/cli.py| changed help text; and the tests gained a fourth regression
baseline, for moist V2, with its numba counterpart. None of these changes
touches a field, function or option the list modifies, and the other source
files in the list are unchanged.

\begin{itemize}
\item \path|ss09/sw_config.py|: fields \path|geometry| (\path|beta_plane|, the
  default, or \path|sphere|), \path|radius| ($6.371\times10^6$ m) and
  \path|rot_rate| ($7.292115\times10^{-5}$ s$^{-1}$); on the sphere, the grid
  $a\,(k/m)$ with $n_y$ odd, the domain fixed at $2a$ (an explicit
  \path|domain_size| or \path|beta| refused), the staggered grid only; the
  geometry arrays, with their flat values on the beta plane.
\item \path|ss09/sw_model.py|: the operators above, through the geometry
  arrays; the collocated layout refuses the sphere, since it is kept frozen to
  reproduce an earlier published version of the model; the output's $y$ described as
  $a\sin\varphi$ on the sphere, with latitude coordinates for the centers and
  faces.
\item \path|ss09/numba_backend.py|: \path|run_day| takes the geometry arrays
  and a switch for the metric term, and each fused loop gains the same factors
  in the same positions as the numpy reference. The shared difference array
  \path|diff| splits in two on the sphere, one for $cu$ and one for $c^2u$,
  which on the plane both equal the current one exactly.
\item \path|ss09/theta_e.py|: the spherical profile, with $\varphi_0$ in
  degrees. The existing profiles are refused on the sphere, since their $y_1$
  and plateau have no meaning there. Seasonal migration of $\varphi_0$ waits
  for the later phase.
\item \path|ss09/hadley_diagnostics.py|: the local Rossby number computed from
  $\partial_y(cu)/f$ in place of $\partial_y u/(\beta y)$, with the call site
  passing $f$ and $c$ in place of $\beta$; jet, cell-center and cell-edge
  positions also reported as latitude. A maximum of $u$ in $y$ is a maximum in
  latitude, since $y$ increases with $\varphi$, so the positions' definitions
  are unchanged.
\item \path|ss09/daily_results.py|: the latitude coordinates and the Rossby
  number's description.
\item \path|ss09/restart_state.py|: the compatibility check compares
  \path|geometry|, \path|radius| and \path|rot_rate|, refusing a beta-plane
  restart for a spherical run.
\item \path|ss09/cli.py|: \path|--geometry|, \path|--radius|,
  \path|--rot-rate|, the spherical forcing type and \path|--lat-0|; on the
  sphere, \path|--beta|, \path|--domain-size|, \path|--y0| and \path|--y1|
  refused.
\item \path|ss09/steady_state.py|: unchanged. Its kinetic-energy metric, the
  mean of $u^2 + v^2$ over the centers, is already an equal-area mean apart
  from counting the two pole half cells as whole cells, and the smoothness
  monitor looks for grid-scale noise in the grid's own coordinate.
\item Tests: new tests of the geometry (antisymmetry, arrays, flat reduction)
  and of the operators (convergence order on regular test fields, the
  dissipation identity, water conservation, the eddy torque, mirror parity over
  multi-day runs); the numba parity suite extended to spherical
  configurations; the regression tests unchanged, which since the merge cover
  four configurations: collocated, staggered, moist V1 and moist V2.
\item Documentation: a section of \path|SCIENCE.md| on the spherical
  formulation, and the new options in \path|README.md| and \path|CLAUDE.md|.
\end{itemize}

\section{Validation, with pass criteria}
\label{sec:valid}

Every criterion is set before its run. Where a criterion tests the model's
operators, an analysis script in this repository's \path|validation/|
directory computes it, independently of those operators.

\paragraph{Code checks, on one machine} (phases 2 and 4).
\begin{enumerate}
\item The flat configuration reproduces the unmodified model bit for bit in the
  four regression configurations, against a reference run, on the same
  machine, of the \path|main| commit the spherical branch starts from. Commit
  \path|6503a62| no longer serves, since the merge changed the moist defaults
  and added the moist V2 configuration.
\item The unit tests above.
\item The numba kernel reproduces the numpy reference bit for bit on the
  sphere, across the existing matrix of physics switches.
\item A hemispherically symmetric forcing keeps exact mirror symmetry over a
  30-day spherical run, and a run continued from a restart file reproduces the
  uninterrupted run bit for bit.
\end{enumerate}

\paragraph{Runs on berimbau} (phase 3), all with the spherical forcing at
$\Delta_y = 50$ K and $\varphi_0 = 0$ unless stated, each plotted and inspected
before any number from it is reported.
\begin{enumerate}[label=(\roman*)]
\item \textbf{Drag-limited angular-momentum identity.} With $v_d = 0$ and the
  vertical advection off, the steady $u$ equation is
  $\bar v\,(f + \zeta) = \epsilon_u u$, so
  $1 - \Ro_{\mathrm{loc}} = \epsilon_u u/(f\bar v)$ at every point, for
  $\Ro_{\mathrm{loc}} = -\zeta/f$. Runs at $\epsilon_u = 10^{-8}$ and
  $10^{-9}$ s$^{-1}$, the second at five drag timescales, 58,000 model days.
  \emph{Pass:} between 5$^\circ$ and the jet, the two sides agree to within the residual of the steady
  $u$ budget, with $\zeta$ computed by the analysis script from $u$. A model
  missing the metric part of the advection would leave a residual of
  $u\tan\varphi/(af)$: 0.057 for $u = 50$ m s$^{-1}$ at 20$^\circ$, and 0.022
  for 20 m s$^{-1}$ at 10$^\circ$.
\item \textbf{Steady budgets} at the production parameters. \emph{Pass:} the
  residuals of the $u$, $v$ and $\theta$ budgets from the final 30-day means,
  relative to each budget's largest term, are no larger than the beta-plane
  twin's of item (iv).
\item \textbf{Timestep and resolution.} The timestep ceiling at $n_y = 325$,
  found by bisection, and a twin at $n_y = 649$, with $\Delta y$ halved.
  \emph{Pass:} the whole-field $\max|\Delta|$ of $u$, $v$ and $\theta$ within
  40$^\circ$ of the equator, each scaled by its field's maximum, is no larger
  than that of the beta-plane twin at the same two spacings. The polar cells
  are reported separately.
\item \textbf{Scaling of the curvature effect.} The sphere beside its
  beta-plane twin: $\beta = 2\Omega/a$, walls at $\pm\pi a/2$, the distance
  from equator to pole, and the same $\theta_E$ as a function of latitude, with
  $\varphi = y/a$. That forcing is the existing profile of Schneider and
  Bordoni (2008) with $y_1 = \pi a/2$ and $y_0 = a\varphi_0$. The two models
  then differ only through curvature: the factors of $c$, the metric term,
  $2\Omega\sin\varphi$ in place of $2\Omega\varphi$, and the spherical form of
  the eddy closure. Each enters at relative order $\varphi^2$, and the square
  of the cell's latitude extent grows with $\Delta_y$, so the relative
  difference between sphere and twin in the jet and cell-edge latitudes should
  scale with the square of the twin's cell-edge latitude, $\varphi_e^2$ in
  radians. Runs at $\Delta_y = 12.5$, 25 and 50 K, with the twin on 508
  intervals ($n_y = 509$, keeping a center on the equator), a spacing 0.18\%
  above the sphere's. \emph{Pass:} that relative difference divided by
  $\varphi_e^2$ is the same across the three rungs to within the numerical
  uncertainty from item (iii), once the next-order term in $\varphi_e^2$ is
  allowed for. A curvature term entering at the wrong order, such as a dropped
  factor of $c$, breaks the proportionality. The runs belong to phase 3; the
  next-order term comes from theory, so this criterion is evaluated in the
  science project of Section~\ref{sec:handoff}. This test
  replaces the earlier plan's ladder of planetary radii at fixed $2\Omega/a$:
  at fixed $\Delta_y$ that ladder narrows the cell in latitude in the same way,
  and it would have required planets other than Earth.
\end{enumerate}

\section{Handoff to the science runs}
\label{sec:handoff}

Phases 2 to 4 belong to this repository. Phase 5, the science runs on
the sphere, belongs to a separate research project in the private repository
\path|spencerahill/fixed-ro-eq-area|. That project takes over once both of the
following hold:
\begin{enumerate}
\item The spherical option is merged into \path|main|, with code checks 1 to 4
  of Section~\ref{sec:valid} passing.
\item Validation runs (i) to (iii) have passed, and the runs of item (iv) are
  complete, with the evaluation of their pass criterion left to the science
  project.
\end{enumerate}
The science project then needs four things from this repository: the commit of
\path|main| that adds the spherical option; the location on berimbau of the
output of validation runs (i) to (iv); the analysis scripts in
\path|validation/| that scored items (i) to (iii); and the measured timestep
ceiling at $n_y = 325$ with the cost per model day on the sphere.

\section{Risks}

\begin{enumerate}
\item \textbf{The polar cells.} The averages between grids and the eddy term
  lose accuracy there, as the polar-cells subsection quantifies. If grid-scale
  noise appears near the poles, in the $v$-smoothness monitor or the
  resolution twin, the first remedy to test is a stronger smoothing confined
  to the outermost cells, a departure from the current numerics tested on its
  own.
\item \textbf{The limiter on $c^2u$.} If the gated eddy term is noisy on the
  sphere, the alternative of limiting $u$ is the first change to test.
\item \textbf{The timestep.} The expectation that the ceiling equals the beta
  plane's at the same $\Delta y$ comes from the equations, and the bisection
  may contradict it.
\item \textbf{The bit-for-bit reduction.} The numba kernel shares one
  difference array between two terms that need different ones on the sphere.
  Splitting it keeps the flat values exactly, but any reordering made while
  splitting would show as a failure of code check 1, and the fix belongs in the
  restructured code.
\end{enumerate}

\section*{References}

\begin{list}{}{\leftmargin=1.5em \itemindent=-1.5em \itemsep=2pt}
\item Adam, O., and N. Paldor, 2009: Global circulation in an axially symmetric
  shallow water model forced by equinoctial differential heating.
  \emph{J. Atmos. Sci.}, \textbf{66}, 1418--1433,
  \url{https://doi.org/10.1175/2008JAS2685.1}.
\item Held, I. M., and A. Y. Hou, 1980: Nonlinear axially symmetric
  circulations in a nearly inviscid atmosphere. \emph{J. Atmos. Sci.},
  \textbf{37}, 515--533.
\item Schneider, T., and S. Bordoni, 2008: Eddy-mediated regime transitions in
  the seasonal cycle of a Hadley circulation and implications for monsoon
  dynamics. \emph{J. Atmos. Sci.}, \textbf{65}, 915--934.
\item Sobel, A. H., and T. Schneider, 2009: Single-layer axisymmetric model for
  a Hadley circulation with parameterized eddy momentum forcing. \emph{J. Adv.
  Model. Earth Syst.}, \textbf{1}, 10,
  \url{https://doi.org/10.3894/JAMES.2009.1.10}.
\item Sobel, A. H., and T. Schneider, 2013: Correction to ``Single-layer
  axisymmetric model for a Hadley circulation with parameterized eddy momentum
  forcing''. \emph{J. Adv. Model. Earth Syst.}, \textbf{5}, 654--657,
  \url{https://doi.org/10.1002/jame.20030}.
\item Xian, P., and R. L. Miller, 2008: Abrupt seasonal migration of the ITCZ
  into the summer hemisphere. \emph{J. Atmos. Sci.}, \textbf{65}, 1878--1895,
  \url{https://doi.org/10.1175/2007JAS2367.1}.
\end{list}

\end{document}
